\section{Introduction}
\label{sec:intro}

The pre-training and fine-tuning paradigm has revolutionized machine learning. 
By starting with a large, expressive foundation model and adapting it to specific tasks, practitioners can deploy highly accurate expert models~\cite{radford2021learning,devlin2018bert}. 
However, deploying a dedicated neural network for every downstream task becomes prohibitively expensive in multi-task, resource-constrained, or edge computing environments. 
Traditional multi-task training or joint fine-tuning offers a potential remedy but incurs high computational costs, requires concurrent access to all training datasets, and is highly prone to gradient conflict and catastrophic interference~\cite{french1999catastrophic,mccloskey1989catastrophic}.

To bypass these limitations, \emph{model merging} has emerged as an active field of study~\cite{seneviratne2024survey,choshen2022fusing}. 
Model merging aims to combine the weights of independent, task-specific expert models (which share a common pre-trained initialization) into a single multi-task network without any retraining or access to the original training data. 
Pioneering work on task arithmetic~\cite{ilharco2022task} showed that subtracting pre-trained base weights from fine-tuned expert weights produces \emph{task vectors}, which can be linearly combined to patch models or inject new capabilities. 
Subsequent methods have developed advanced techniques to handle parameter interference, such as TIES-Merging~\cite{yadav2023ties}, Fisher-weighted averaging~\cite{matena2022merging}, and permutation-based weight alignment~\cite{ainsworth2023git,zhang2023zipit}.

Despite these advances, standard merging recipes often rely on manually-tuned global coefficients to scale task vectors. 
To address this, recent works have proposed \emph{adaptive model merging} (e.g., AdaMerging~\cite{yang2024adamerging}), where the merging coefficients are optimized dynamically on a small, unlabeled test-time calibration stream by minimizing an unsupervised objective, such as prediction entropy. 
While adaptive merging can outperform static baselines when operating on fully-converged, high-capacity foundation models, we empirically demonstrate that in resource-constrained or low-fidelity edge regimes, unconstrained adaptation actually degrades performance compared to a simple uniform weight blend. In fact, in our challenging setting, the static Uniform Task Arithmetic baseline of 30.41\% remains the superior model, as transductive noise on compact calibration streams ($N=256$) completely overwhelms fine-grained, unregularized updates. This surprising finding highlights a fundamental, unexamined question about the \textbf{structural resolution of parameter optimization}:
\begin{center}
\emph{At what level of physical granularity should we define and optimize merging coefficients?}
\end{center}

Prior works are highly fragmented across this structural spectrum. 
While task arithmetic~\cite{ilharco2022task} operates at the coarsest level of a single global scalar per task, AdaMerging~\cite{yang2024adamerging} optimizes layer-wise coefficients (one scalar per Transformer layer per task), and SplineMerge/PolyMerge~\cite{jung2025poly} parameterize layer coefficients as continuous low-degree splines. 
At the other extreme, one could theoretically optimize an individual coefficient for every single weight tensor and bias vector in the model (tensor-wise merging), or even parameter-wise. 
Intuitively, increasing the granularity of merging coefficients increases the optimization capacity, allowing the model to route features and resolve parameter interference with surgical precision. 
However, as we increase the parameter count from a single global scalar to hundreds of tensor-wise coefficients, we run headfirst into a classic machine learning conflict: \textbf{the Generalization-Granularity Trade-off}.

Optimizing high-dimensional merging coefficients on small test-time calibration streams (e.g., $N=256$ samples) is highly susceptible to \emph{transductive overfitting}~\cite{jin2026regcal,mueller2023analyzing}. 
In this regime, the optimizer may find extreme, unphysical coefficient values that minimize the surrogate unsupervised loss (entropy) on the calibration batch but destroy the model's underlying representations, leading to a catastrophic drop in true multi-task generalization accuracy. 
Thus, understanding and mapping this trade-off is critical for designing stable, deployment-ready adaptive merging systems.

\begin{figure}[t]
\centering
\includegraphics[width=\columnwidth]{granularity_tradeoff.png}
\caption{\textbf{The Generalization-Granularity Trade-off in Model Merging.} Main-task generalization performance across five levels of structural granularity. On compact test-time calibration streams ($N=256$), increasing parameter granularity leads to progressive transductive overfitting, degrading global multi-task generalization. Black-box zero-order 1+1 ES maintains much greater robustness than unconstrained Adam gradient descent, a phenomenon we attribute to a combination of isotropic search boundaries and optimization sluggishness.}
\label{fig:tradeoff_curve}
\end{figure}

As committed \textbf{Empiricists}, we believe that such fundamental trade-offs cannot be resolved by speculative hand-waving or minor heuristic adjustments. 
They demand systematic, large-scale experimentation, exhaustive sweeps across diverse hyperparameter spaces, and rigorous statistical validation. 
To this end, we introduce \textbf{GranMerge}, a unified empirical framework designed to dissect the structural boundaries of adaptive weight blending. 
We evaluate GranMerge across five nested levels of parameter resolution (from Global to Tensor-wise) across four visual tasks (MNIST, FashionMNIST, CIFAR-10, SVHN) using a 12-layer Vision Transformer backbone, running exhaustive multi-seed sweeps under two distinct optimization families: first-order gradient descent (Adam) and zero-order evolution (1+1 ES).

Our extensive empirical investigation yields several major insights:
\begin{enumerate}
    \item \textbf{The Degradation of Generalization with Granularity:} We show that as parameter resolution increases, unconstrained test-time adaptation suffers from severe transductive overfitting. Rather than resolving task interference, fine-grained optimization (Level 5 Tensor-wise) exploits local prediction entropy, degrading global multi-task generalization compared to simple global optimization (Level 1) or uniform baselines.
    \item \textbf{First-order vs. Zero-order Overfitting Dynamics:} We show that Adam gradient descent is highly vulnerable to transductive overfitting, as unconstrained gradients rapidly exploit high-dimensional parameters to minimize local entropy. In contrast, black-box zero-order 1+1 Evolution Strategies (ES) maintain stable generalization, which we empirically deconstruct as a combination of isotropic walk constraints and optimization sluggishness (underfitting) under high dimensions.
    \item \textbf{The Limitations of Simple Regularizers:} We introduce and benchmark Elastic Spatial Regularization (ESR) and depth-wise Total Variation (TV) smoothness penalties. We show that while these L2 constraints successfully stabilize zero-order ES (improving Level 5 ES from 29.43\% to 30.17\% and Adam from 26.91\% to 28.51\%), they are insufficient to arrest the extreme, chaotic overfitting of first-order Adam. This highlights a key limitation of first-order test-time adaptation and underscores the need for stronger structural constraints (such as low-degree splines) in gradient-based blending.
\end{enumerate}

The remainder of this paper is structured as follows: Section~\ref{sec:related} reviews related work; Section~\ref{sec:method} formalizes the GranMerge framework, its granularities, and regularizers; Section~\ref{sec:experiments} presents our extensive empirical results, trade-off curves, limitations, and ablation studies; and Section~\ref{sec:conclusion} concludes with actionable insights for practitioners.
